\subsection*{Efecto Ramsauer--Townsend}

Las mediciones de las Figuras 5 y 6 muestran una respuesta no
monótona del voltaje de ánodo $V_a$. Después de una región de crecimiento,
$V_a$ alcanza un máximo amplio y posteriormente disminuye, mientras que el
voltaje de rejilla-blindaje $V_{rb}$ mantiene un crecimiento aproximadamente
monótono. Puesto que $V_a$ es proporcional a la corriente de ánodo, este
máximo indica un intervalo en el que una mayor fracción de los electrones
emitidos alcanza el ánodo. Este comportamiento es compatible con la disminución de la sección eficaz
de dispersión asociada con el efecto Ramsauer--Townsend en xenón
\cite{Kukolich1968}.\\

La adquisición automatizada reproduce la tendencia general de crecimiento y
disminución de $V_a$, así como el aumento de $V_{rb}$; sin embargo, presenta
una dispersión considerable entre puntos consecutivos. La
diferencia sugiere que, durante la adquisición automatizada, el tiempo de
estabilización después de cada cambio de $V_c$ pudo ser insuficiente o que
existieron fluctuaciones en el canal de medición de $V_a$. Bajo las
condiciones empleadas, las curvas manuales resultaron más adecuadas para
identificar la posición y la forma general del máximo.\\

Para mejorar la comparación entre ambos métodos sería conveniente realizar
varios barridos para cada valor de $V_f$, introducir un tiempo de espera fijo
después de modificar $V_c$ y promediar varias lecturas en cada punto.



\subsection*{Técnica de Townsend pulsada}

En la Figura \ref{v_contra_Tausen}, los datos obtenidos para diferentes presiones se superponen
aproximadamente cuando la velocidad de deriva se representa en función del
campo eléctrico reducido $E/N$. Esta coincidencia indica que, dentro del intervalo estudiado, el transporte
electrónico está determinado principalmente por $E/N$ y no de manera
independiente por el campo eléctrico o la presión, como se espera en los
experimentos de enjambres electrónicos \cite{GonzalezMagana2019}.\\

La velocidad de deriva aumenta inicialmente hasta alcanzar un máximo cercano
a $E/N\approx 5$--$6~\mathrm{Td}$. Después disminuye hasta un mínimo amplio
alrededor de $E/N\approx 10$--$11~\mathrm{Td}$, antes de volver a crecer para
campos reducidos mayores. La región descendente satisface
$
    d v_d/d(E/N)<0
$, 
por lo que constituye evidencia de conductividad diferencial negativa. Este
comportamiento concuerda cualitativamente con lo reportado para mezclas de
argón con pequeñas concentraciones de nitrógeno \cite{Jovanovic2009}. En esta
mezcla, la forma de la distribución de energía de los electrones está
determinada por la competencia entre el mínimo Ramsauer--Townsend del argón
y los canales de excitación del nitrógeno. Como resultado, un aumento del
campo reducido no produce necesariamente un incremento inmediato de la
velocidad media de los electrones.\\

Para valores elevados de $E/N$, la velocidad vuelve a aumentar de manera
pronunciada. En este régimen, la energía adquirida por los electrones entre
colisiones se vuelve suficiente para superar la región asociada con la
conductividad diferencial negativa. No obstante, esta parte de la curva fue
medida principalmente a las presiones más bajas, debido a que en las demás
condiciones aparecieron descargas entre los electrodos y las paredes de la
cámara. En consecuencia, la superposición entre distintas presiones está
bien sustentada en la región de campos bajos e intermedios, pero no puede
evaluarse con la misma precisión en el régimen de campos altos.\\

En conclusión, ambas experiencias permitieron estudiar el transporte electrónico desde perspectivas complementarias. El experimento del efecto Ramsauer--Townsend puso de manifiesto la influencia de la sección eficaz de dispersión sobre el movimiento de electrones individuales, mientras que la técnica de Townsend pulsada permitió caracterizar el comportamiento colectivo de un enjambre electrónico mediante la determinación de su velocidad de deriva. Los resultados obtenidos fueron consistentes con la teoría revisada durante el curso y permitieron comprender la relación existente entre los procesos microscópicos de colisión y las propiedades macroscópicas del transporte de carga en gases.